\documentclass[11pt]{article}\usepackage[margin=1in]{geometry}\usepackage{amsmath,amsthm}\newtheorem{theorem}{Theorem}\newtheorem{remark}{Remark}
\title{The Link-Lift Spillover Obstruction for a Single-Star Repair}\author{}\date{21 July 2026}\begin{document}\maketitle
\begin{abstract}We lift a centered six-plaquette repair to the links of the open $2^4$ block and compute the least external curvature it must create.\end{abstract}
Let $A$ map noncentral link variations to the five centered incident-curvature modes and let $B$ map them to the other 210 plaquettes. Define
\[H=R^*R,\qquad R y=\mathop{\rm argmin}_{Ax=y}\|Bx\|.
\]
\begin{theorem}For each of the eight six-plaquette interior links, the spectrum of $H$ is
\[\{0,\ 1.121212\ldots,\ 1.121212\ldots,\ 1.635983\ldots,\ 1.635983\ldots\}.
\]
Thus only one centered mode has a zero-spill lift; four modes necessarily export curvature.
\end{theorem}
\begin{proof}Choose an orthonormal basis of the centered star. A particular right inverse is $L_0=A^*(AA^*)^{-1}$. Varying it by $\ker A$ shows that the minimum external field is the orthogonal residual of $BL_0$ modulo $B\ker A$. The displayed finite matrix $R^*R$ is diagonalized by the executable certificate; all eight incidence stars are congruent.\end{proof}
After action drop, fifteen-dimensional Jacobian, Haar loss, and local comparison loss are combined, the worst-mode single-star repair misses the required free energy even assuming zero external curvature before repair.
\begin{remark}The ideal Stage 43 repair therefore does not lift locally. A coordinated polymer repair must reduce the effective spillover spectrum.\end{remark}\end{document}
